%---------------------------------------------------------%
%   CHAPTER 2: TECHNICAL AND FORMATTING SPECIFICATIONS     %
%---------------------------------------------------------%
\chapter{Technical and Formatting Specifications}

This chapter summarises the formatting rules from the UMP Thesis Guide 2019 \citep{umpthesis2019}. Each subsection ends with a \textbf{LaTeX note} that shows how the rule is applied when using \texttt{umpthesis.cls}.

\section{Paper Specifications}

Table \ref{tab:paperspec} lists the paper requirements.

\begin{table}[htbp]
\caption{Paper specifications}
\label{tab:paperspec}
\begin{tabularx}{\textwidth}{@{} l X @{}}
\toprule
Element & Specification \\
\midrule
Size & A4 (210 $\times$ 297 mm) for both soft copy and hard copy. \\
Quality & Original: simile paper, minimum 80 gsm. Duplicates: good quality photocopy or offset paper. Isograph and carbon copies are not acceptable. \\
Colour & White paper only. \\
Printing & One side of the paper only. The original must be printed with a laser printer. \\
Binding & Final submission: three hard-bound copies (one original and two duplicates). Master's cover is black; PhD cover is maroon. For undergraduate PSM reports, cover colour is at the faculty's discretion. \\
\bottomrule
\end{tabularx}
\end{table}

\textbf{LaTeX note.} The class loads \texttt{a4paper} and \texttt{oneside}. Print one-sided from the PDF viewer.

\section{Typeface or Font}

All text, tables, and graphs must be typeset, not handwritten. Use \textbf{Times New Roman 12 pt} throughout the report, including headings and page numbers. Table entries may use Times New Roman 11 pt or smaller, as long as they remain legible. Equations must be produced with an equation editor, including subscripts and superscripts.

\textbf{LaTeX note.} The class loads \texttt{mathptmx} (the standard Times New Roman substitute in \LaTeX{}) at \texttt{12pt}:
\begin{verbatim}
\documentclass[english]{umpthesis}   % 12pt Times-like font
\end{verbatim}

\section{Margin Specifications}

All text and graphics must stay within the margins in Table \ref{tab:margins}.

\begin{table}[htbp]
\caption{Margin specifications for A4 paper}
\label{tab:margins}
\begin{tabularx}{\textwidth}{@{} l >{\centering\arraybackslash}X >{\centering\arraybackslash}X @{}}
\toprule
Element & Portrait & Landscape \\
\midrule
Top edge & 25 mm & 35 mm \\
Right side & 25 mm & 25 mm \\
Left side & 35 mm & 25 mm \\
Bottom edge & 25 mm & 25 mm \\
\bottomrule
\end{tabularx}
\end{table}

\textbf{LaTeX note.} The class itself does not set the body margins, so the main file must set them explicitly with \texttt{geometry}:
\begin{verbatim}
\geometry{top=25mm,right=25mm,left=35mm,bottom=25mm,includefoot}
\end{verbatim}
For landscape pages, use \texttt{lscape} or \texttt{pdflscape} and set the landscape margins with \texttt{\string\newgeometry} (see Chapter \ref{chap:latex}).

\section{Line Spacing}

Table \ref{tab:spacing} summarises the spacing rules.

\begin{table}[htbp]
\caption{Spacing formatting}
\label{tab:spacing}
\begin{tabularx}{\textwidth}{@{} l X @{}}
\toprule
Line spacing & Applies to \\
\midrule
25 mm & Spacing between the upper edge of the page and the chapter number. \\
1.5-line spacing & Between the chapter number and chapter title; between the title and the first line of text; between the last line of text and a subsection title; between a subsection title and the first line of text; between paragraphs; between the last line of text and a table or figure; list of tables, figures, symbols, and abbreviations; general body text. \\
Single (1.0) & Acknowledgements; abstract and its translation; long title entries in headings; long tables and table entries; long quotations, notes, and multiline captions; footnotes; bibliographic entries/references. \\
\bottomrule
\end{tabularx}
\end{table}

\textbf{LaTeX note.} The class sets the body to approximately 1.5 spacing when \texttt{\string\mainmatter} begins:
\begin{verbatim}
\linespread{1.435}          % set by umpthesis.cls
\end{verbatim}
Acknowledgements, abstract, and dedication environments use \texttt{\string\singlespacing} automatically. Known spacing deviations are listed in Chapter \ref{chap:checklist}.

\section{Paragraphing}

All paragraphs must have the first line indented \textbf{12.7 mm} from the left margin. The last paragraph on a page must have at least two lines; otherwise it must begin on the next page. The last word on a page must not be split with a hyphen.

\textbf{LaTeX note.} The class sets the indent and penalties:
\begin{verbatim}
\setlength{\parindent}{0.5in}   % 12.7 mm, set by umpthesis.cls
\widowpenalty10000
\clubpenalty10000
\end{verbatim}

\section{Pagination and Numbering}

Page numbers use the same Times New Roman font as the text, are centred, and sit within \textbf{20 mm} of the bottom edge. Preliminaries are numbered in lowercase Roman numerals (i, ii, iii). The text uses Arabic numerals (1, 2, 3). The first page of the preliminaries is counted but its number is not printed. Letter suffixes (10a, 10b) are not allowed.

\textbf{LaTeX note.} The class uses \texttt{\string\frontmatter} and \texttt{\string\mainmatter} to switch numbering:
\begin{verbatim}
\frontmatter   % roman page numbers for preliminaries
\mainmatter    % arabic page numbers for the text
\end{verbatim}
Known issue: the class currently hides the page number on the first page of every chapter (see Chapter \ref{chap:checklist}).

\section{Front Cover, Spine, and Title Page}

\subsection{Front Cover}

The front cover uses \textbf{gold coloured CAPITAL LETTERS, not bold}, in Times New Roman \textbf{20 pt}, in this order:
\begin{enumerate}
    \item Title of the report (50 mm from the top edge of the cover);
    \item Full name of the student (as in identity card or passport);
    \item Name of the degree (e.g., \textsc{Bachelor of Computer Science});
    \item University name (\textsc{Universiti Malaysia Pahang}, 50 mm from the bottom edge).
\end{enumerate}

\textbf{LaTeX note.} Use \texttt{\string\umpCoverPage}. The class renders this page at 18 pt and does not apply the gold colour (see Chapter \ref{chap:checklist}).

\subsection{Spine}

The spine shows, in gold capital letters, Times New Roman 18 pt: student name, abbreviated degree, year of award, and \textsc{UMP}.

\textbf{LaTeX note.} The class does not currently provide a spine command.

\subsection{Title Page}

The title page uses \textbf{single spacing and CAPITAL LETTERS}, in this order:
\begin{enumerate}
    \item Title of the report (same as the cover, 50 mm from the top edge);
    \item Full name of the student;
    \item Purpose of submission (e.g., \emph{Report submitted in partial fulfilment of the requirements for the award of the degree of Bachelor of \dots});
    \item Faculty/Institute/Centre (Title Case);
    \item \textsc{Universiti Malaysia Pahang} (uppercase);
    \item Month and year of award (50 mm from the bottom edge).
\end{enumerate}

\textbf{LaTeX note.} Use \texttt{\string\maketitle}. Known issue: the class title page uses 1.5 spacing instead of single spacing (see Chapter \ref{chap:checklist}).

\section{Headings and Subheadings}

Table \ref{tab:headings} gives the formatting style for each heading level.

\begin{table}[htbp]
\caption{Formatting of headings and subheadings}
\label{tab:headings}
\begin{tabularx}{\textwidth}{@{} l X @{}}
\toprule
Heading level & Formatting \\
\midrule
Chapter title (Heading 1) & Centred, boldface, uppercase. Every chapter starts on a new page. \\
Section title (Heading 2, UMP Level 2) & Left aligned, boldface, Title Case. Paragraph begins below, indented as a normal paragraph with 1.5 spacing. Numbered 2.1. Included in the Table of Contents. \\
Subsection title (Heading 3, UMP Level 3) & Left aligned, boldface, Title Case. Paragraph begins below, indented with 1.5 spacing. Numbered 2.1.1. Included in the Table of Contents. \\
Level 4 (UMP Level 4) & Left aligned, boldface, Title Case. Numbered 2.1.1.1. Not included in the Table of Contents. \\
Level 5 (UMP Level 5) & Left aligned, not numbered, italicized, Title Case. Not included in the Table of Contents. \\
Heading 7 & Roman numbering (i, ii, iii). Use for objectives and scope. Not in Table of Contents. \\
Heading 8 & Arabic numbering (1, 2, 3). Not in Table of Contents. \\
Heading 9 & Bullet numbering. Not in Table of Contents. \\
\bottomrule
\end{tabularx}
\end{table}

\textbf{LaTeX note.} The class maps these levels as follows:
\begin{verbatim}
\chapter{...}        % Heading 1
\section{...}        % Heading 2 (2.1)
\subsection{...}     % Heading 3 (2.1.1)
\subsubsection{...}  % Heading 4 (2.1.1.1)
\paragraph{...}      % Heading 5 (see known issue)
\begin{enumerate}[label=(\roman*)] ... \end{enumerate}  % Heading 7
\begin{enumerate} ... \end{enumerate}                   % Heading 8
\begin{itemize} ... \end{itemize}                       % Heading 9
\end{verbatim}

\section{Tables}

Number tables with Arabic numerals by chapter: the first table in Chapter 3 is \texttt{Table 3.1}. Every table needs a short caption. The caption is always placed \textbf{above} the table, even for landscape tables. If a table is copied or adapted from a source, cite the source at the bottom:
\begin{verbatim}
Source: Author (Year).
\end{verbatim}

A table must contain only three horizontal lines and no vertical lines: an opening line below the title, a line that closes the column headings, and a closing line at the bottom. Every column must have a heading. Continued tables must repeat the column headings and show \emph{continued} with the table number on subsequent pages.

\textbf{LaTeX note.} Use \texttt{booktabs} for the three-line style and place \texttt{\string\caption} before the tabular environment:
\begin{verbatim}
\begin{table}[htbp]
\caption{Table caption goes above}
\label{tab:example}
\begin{tabularx}{\textwidth}{@{} l X @{}}
\toprule
Heading A & Heading B \\
\midrule
Value & Description \\
\bottomrule
\end{tabularx}
\end{table}
\end{verbatim}

\section{Figures}

Figures (maps, charts, graphs, drawings, diagrams, photographs) must be clear and high quality. Number figures by chapter: \texttt{Figure 4.2} is the second figure in Chapter 4. Text inside a figure may use any neat, readable typeface, but the \textbf{figure number and caption} must use the same typeface as the body text. The caption is placed \textbf{below} the figure and is left aligned. Figures are centred between the left and right margins. If a figure is in landscape orientation, the top of the figure is placed along the bound edge.

\textbf{LaTeX note.} The class configures the caption package for left-aligned captions:
\begin{verbatim}
\begin{figure}[htbp]
\centering
\includegraphics[width=0.8\textwidth]{figure/example.png}
\caption{Caption goes below the figure}
\label{fig:example}
\end{figure}
\end{verbatim}

\section{Equations}

Equations are centred. Equation numbers use Arabic numerals in parentheses on the right margin and are numbered by chapter: \texttt{Eq. 4.2} is the second equation in Chapter 4. Break long equations before a convenient plus or minus sign, or just ahead of an equal sign.

\textbf{LaTeX note.} The memoir base class numbers equations by chapter automatically:
\begin{verbatim}
\begin{equation}
E = mc^{2}
\end{equation}
\end{verbatim}

\section{Citations and References}

In-text citations follow the APA 6th author--date style. Every source cited in the text must appear in the References list, and every entry in the list must be cited in the text. The References list starts on a new page with the heading \textsc{References} centred, 25 mm from the top. References are not listed at the end of each chapter; they are listed once at the end of the final chapter, in alphabetical and chronological order, with hanging indentation.

\textbf{LaTeX note.} The class does not provide an APA style file. The current projects use \texttt{natbib} with \texttt{plainnat} (author--year, close to APA but not exact):
\begin{verbatim}
\usepackage{natbib}
...
\citet{key}   % Author (Year)
\citep{key}   % (Author, Year)
...
\bibliographystyle{plainnat}
\bibliography{bib/mybiblio}
\end{verbatim}
For strict APA 6th output, use a dedicated APA bibliography style or \texttt{biblatex-apa}. This is a known gap (see Chapter \ref{chap:checklist}).

\section{Appendices and List of Publications}

Appendices hold material that is relevant but too lengthy for the text (raw data, questionnaires, program listings, maps, etc.). They are paginated consecutively with the main text. If there are more than five appendices, include a \emph{List of Appendices} in the Table of Contents. Publications resulting from the study during candidature should be listed clearly in a \emph{List of Publications}.

\textbf{LaTeX note.} Use \texttt{\string\appendix} followed by \texttt{\string\chapter\string{...\string}}. The class does not currently define commands for List of Appendices, List of Publications, or List of Symbols/Abbreviations; define them manually or add them to the class.
